\section{Discussion}

This chapter proposed a property-driven training and enforcement for safety-critical systems that use ANN-based modules to address the challenges associated with designing and ensuring safe behaviour in complex cyber-physical systems. The proposed approach uses a combination of compositional design, property-driven training, and runtime enforcement to ensure that the system satisfies the desired safety properties.

We begin by introducing the safety-critical system considered as a case study -  an Autonomous Vehicle system. The system is designed to be modular or compositional, with each module responsible for a specific task such as Car Following (CF), Lane Change Decision (LCD), and Lane Change Execution (LCE). We then discuss safety in the autonoumous vehicle system and the challenges associated with ensuring safety in such systems. We propose a compositional property-driven training and enforcement approach to address these challenges.

We continue by presenting the novel two-step property-driven training approach. The first step uses a convex optimisation process with the desired safety properties (provided as templates with unknown coeffients) used as constraints during the optimisation to mine properties from data. The optimisation stops when a desired threshold of safety property satisfaction is achieved. The second step uses the mined properties to produce a new dataset from the training dataset such that the new dataset satisfies the property. This dataset can then be used to train the ANN-based modules. The model agnostic nature of the proposed approach makes it suitable for compositional ANN models where each model may have its own set of safety properties. 

Next, we present the Runtime Enforcement mechanism used in this work and the process for enforcer synthesis. The runtime enforcement mechanism uses a set of defined properties to monitor the system's behaviour and take corrective actions if the system violates the desired safety properties. This is followed by a formal definition of the properties used to define safe behaviour in the AV systemm, where we illustrate both system level and module level properties. We follow this with a discussion on how the compositional LCD module is trained using the proposed property-driven training approach. We show the training of both monolithic and compositional models and compare the results obtained from both training methodologies.

Finally, we discuss the simulation set-up used to simulate the AV scenario and the results obtained from the simulation. We show that the compositional models trained with property-driven methodologies outperform monolithic models in enforcing safety-critical behavior in autonomous vehicle systems. Specifically, the compositional approach resulted in significantly fewer runtime faults compared to monolithic models across multiple simulations. Additionally, the runtime enforcers effectively corrected safety violations, ensuring adherence to predefined policies in all scenarios tested. These findings affirm the suitability of property-driven compositional designs for complex safety-critical systems. 
